\section{Related Work}

Eye tracking has long been used as behavioural evidence for internal user
states. Pupil dilation is a classical correlate of processing load
\cite{kahneman1966pupil}, and fixation, saccade, blink and pupil statistics
computed over short windows remain the dominant input to models of workload,
affect and engagement. Our own prior work illustrates both the appeal and the
limits of this pipeline: handcrafted ocular features supported coarse
affective distinctions \cite{bozak2024emotion}, yet a systematic evaluation
across several prediction targets showed that task type is recognised far more
reliably than graded workload, and that reported accuracy depends strongly on
how the prediction target and the validation split are defined
\cite{bozak2026workload}. Feature engineering encodes useful domain knowledge,
but the resulting performance is bounded by what a fixed feature set captures.

Learned representations offer an alternative. GazeMAE pretrains temporal
autoencoders on raw position and velocity signals and shows that the resulting
gaze-specific embeddings support downstream classification through a simple
linear head \cite{bautista2020gazemae}. In parallel, general-purpose
time-series foundation models such as MOMENT provide a frozen multichannel
encoder intended to transfer across domains without task-specific pretraining
\cite{goswami2024moment}. Both are presented as replacements for
hand-designed features, and both are typically evaluated on curated
benchmarks rather than on modest, subject-dependent recordings of spontaneous
behaviour. How much of that benefit survives when such encoders are applied
out of the box in this setting remains open.

Evaluation design is a further concern. Windows drawn from the same person are
highly dependent, so random window splits and preprocessing fitted on the full
dataset inflate reported performance, an effect documented in detail for other
biosignal domains \cite{rosenblatt2024leakage}. Comparisons across
representations are interpretable only when every candidate is evaluated on
identical windows and labels under a subject-disjoint protocol with fold-local
preprocessing.

We therefore extend our earlier feature-based analyses to learned
representations. This preliminary study is an empirical comparison rather than
a new encoder or fine-tuning method: raw signals, handcrafted features,
GazeMAE and MOMENT embeddings are evaluated on the same windows, the same
engagement target and the same leave-one-subject-out folds.
